\documentclass[10pt,a4paper]{amsart}
\usepackage{amsmath,amssymb,amsthm,graphicx}
% Retained preamble for the Stacks Project chapter {CH} (shared project preamble).
% The chapter file pinned-original-source.tex only carries \input{{preamble}} on its first line,
% so no preamble text lives inside the chapter itself.  The Stacks Project uses one single shared
% preamble for every chapter: the snapshot stages it three times (sources/stacks-chow/preamble.tex,
% sources/stacks-etale-cohomology/preamble.tex and focused-seeds/AGP/sources/stacks-snapshot-20261011/
% preamble.tex) and those copies are identical in size.  The block below is copied mechanically,
% line for line, from sources/stacks-chow/preamble.tex lines 14-18 and 154-217: the Xy-pic setup and
% the native theorem declarations together with the mathematical macros.  Nothing was added, removed
% or redefined; the wrapper class and the amsmath/amssymb/amsthm/graphicx base packages are supplied
% by the assembler.
% For commutative diagrams we use Xy-pic
% Chapter: stacks-spaces-more-morphisms.tex -- More on Morphisms of Spaces.
\usepackage[all]{xy}

% We use 2cell for 2-commutative diagrams.
\xyoption{2cell}

% Theorem environments.
%
\theoremstyle{plain}
\newtheorem{theorem}[subsection]{Theorem}
\newtheorem{proposition}[subsection]{Proposition}
\newtheorem{lemma}[subsection]{Lemma}

\theoremstyle{definition}
\newtheorem{definition}[subsection]{Definition}
\newtheorem{example}[subsection]{Example}
\newtheorem{exercise}[subsection]{Exercise}
\newtheorem{situation}[subsection]{Situation}

\theoremstyle{remark}
\newtheorem{remark}[subsection]{Remark}
\newtheorem{remarks}[subsection]{Remarks}

\numberwithin{equation}{subsection}

% Macros
%
\def\lim{\mathop{\mathrm{lim}}\nolimits}
\def\colim{\mathop{\mathrm{colim}}\nolimits}
\def\Spec{\mathop{\mathrm{Spec}}}
\def\Hom{\mathop{\mathrm{Hom}}\nolimits}
\def\Ext{\mathop{\mathrm{Ext}}\nolimits}
\def\SheafHom{\mathop{\mathcal{H}\!\mathit{om}}\nolimits}
\def\SheafExt{\mathop{\mathcal{E}\!\mathit{xt}}\nolimits}
\def\Sch{\mathit{Sch}}
\def\Mor{\mathop{\mathrm{Mor}}\nolimits}
\def\Ob{\mathop{\mathrm{Ob}}\nolimits}
\def\Sh{\mathop{\mathit{Sh}}\nolimits}
\def\NL{\mathop{N\!L}\nolimits}
\def\CH{\mathop{\mathrm{CH}}\nolimits}
\def\proetale{{pro\text{-}\acute{e}tale}}
\def\etale{{\acute{e}tale}}
\def\QCoh{\mathit{QCoh}}
\def\Ker{\mathop{\mathrm{Ker}}}
\def\Im{\mathop{\mathrm{Im}}}
\def\Coker{\mathop{\mathrm{Coker}}}
\def\Coim{\mathop{\mathrm{Coim}}}

% Boxtimes
%
\DeclareMathSymbol{\boxtimes}{\mathbin}{AMSa}{"02}

%
% Macros for moduli stacks/spaces
%
\def\QCohstack{\mathcal{QC}\!\mathit{oh}}
\def\Cohstack{\mathcal{C}\!\mathit{oh}}
\def\Spacesstack{\mathcal{S}\!\mathit{paces}}
\def\Quotfunctor{\mathrm{Quot}}
\def\Hilbfunctor{\mathrm{Hilb}}
\def\Curvesstack{\mathcal{C}\!\mathit{urves}}
\def\Polarizedstack{\mathcal{P}\!\mathit{olarized}}
\def\Complexesstack{\mathcal{C}\!\mathit{omplexes}}
% \Pic is the operator that assigns to X its picard group, usage \Pic(X)
% \Picardstack_{X/B} denotes the Picard stack of X over B
% \Picardfunctor_{X/B} denotes the Picard functor of X over B
\def\Pic{\mathop{\mathrm{Pic}}\nolimits}
\def\Picardstack{\mathcal{P}\!\mathit{ic}}
\def\Picardfunctor{\mathrm{Pic}}
\def\Deformationcategory{\mathcal{D}\!\mathit{ef}}
\usepackage{cleveref}
\crefname{originallemma}{original lemma}{original lemmas}
\Crefname{originallemma}{Original lemma}{Original lemmas}
\makeatletter
\newlabel{lemma-where-unramified}{{05X8}{source}{}{sourcepointer-1}{}}
\newlabel{lemma-where-unramified@cref}{{[originallemma][0][]05X8}{[1][0][]source}{}{}{}}
\makeatother
\crefname{originalreference}{original reference}{original references}
\Crefname{originalreference}{Original reference}{Original references}
\makeatletter
\newlabel{spaces-morphisms-lemma-diagonal-unramified-morphism}{{05W1}{source}{}{sourcepointer-2}{}}
\newlabel{spaces-morphisms-lemma-diagonal-unramified-morphism@cref}{{[originalreference][0][]05W1}{[1][0][]source}{}{}{}}
\makeatother
\crefname{originalreference}{original reference}{original references}
\Crefname{originalreference}{Original reference}{Original references}
\makeatletter
\newlabel{spaces-morphisms-lemma-universally-injective}{{040X}{source}{}{sourcepointer-3}{}}
\newlabel{spaces-morphisms-lemma-universally-injective@cref}{{[originalreference][0][]040X}{[1][0][]source}{}{}{}}
\makeatother
\crefname{originalreference}{original reference}{original references}
\Crefname{originalreference}{Original reference}{Original references}
\makeatletter
\newlabel{spaces-properties-lemma-factor-through-open-subspace}{{03IE}{source}{}{sourcepointer-4}{}}
\newlabel{spaces-properties-lemma-factor-through-open-subspace@cref}{{[originalreference][0][]03IE}{[1][0][]source}{}{}{}}
\makeatother
\crefname{originalreference}{original reference}{original references}
\Crefname{originalreference}{Original reference}{Original references}
\makeatletter
\newlabel{spaces-properties-lemma-open-subspaces}{{03BZ}{source}{}{sourcepointer-5}{}}
\newlabel{spaces-properties-lemma-open-subspaces@cref}{{[originalreference][0][]03BZ}{[1][0][]source}{}{}{}}
\makeatother
\crefname{originalreference}{original reference}{original references}
\Crefname{originalreference}{Original reference}{Original references}
\makeatletter
\newlabel{spaces-properties-lemma-points-cartesian}{{03H4}{source}{}{sourcepointer-6}{}}
\newlabel{spaces-properties-lemma-points-cartesian@cref}{{[originalreference][0][]03H4}{[1][0][]source}{}{}{}}
\makeatother
\title{The locus where a map is unramified is open inside the locus where it is universally injective}
\author{The Stacks Project Authors}
\date{stacks-project:spaces-more-morphisms.tex:local-snapshot-20261011; GNU FDL 1.2 or later}
\begin{document}
\maketitle
\noindent\textit{Editorial note.} Counted span is the author's own proof environment, source lines 11586-11639, immediately following the statement at 11565-11584.\par\medskip
\noindent\textit{Editorial note.} Five cross-chapter pointers (Morphisms of Spaces, Algebraic Spaces) plus a same-chapter companion are recorded in assembly.external\_references.\par\medskip
\noindent\textit{Editorial note.} The proof produces an open subspace and verifies the factorisation property; both steps are inside the counted span.\par\medskip
\noindent\textit{Allowed standard prerequisites (editorial).}\begin{itemize}
\item The unramified locus of a morphism. chapter lemma lemma-where-unramified, cited in the proof. Supplies the corresponding statement for unramifiedness which is adapted here to the universally injective setting.
\item Criterion for universal injectivity. Morphisms of Spaces, lemma spaces-morphisms-lemma-universally-injective, cited in the proof. Turns universal injectivity into a condition on the diagonal, which is what the closedness hypotheses are used to test.
\item The diagonal of an unramified morphism. Morphisms of Spaces, lemma spaces-morphisms-lemma-diagonal-unramified-morphism, cited in the proof. Relates the unramified locus to openness of the diagonal.
\item Open subspaces and factorisation through them. Algebraic Spaces, lemmas spaces-properties-lemma-open-subspaces, spaces-properties-lemma-factor-through-open-subspace and spaces-properties-lemma-points-cartesian, cited in the proof. Convert the set-theoretic conditions obtained in the proof into an open subspace with the required factorisation property.
\end{itemize}\medskip
\ifdefined\hypertarget\hypertarget{sourcepointer-1}{}\fi
\noindent\textit{Source reference.} Original lemma 05X8: sources/stacks-spaces-more-morphisms/pinned-original-source.tex, source lines 11529--11529. Navigation pointer only: the counted proof cites this statement by \textbackslash{}ref as a standing input; the deduction the counted proof has to make is carried out inside the counted span, and no part of that deduction is delegated to this pointer.\par\medskip
\ifdefined\hypertarget\hypertarget{sourcepointer-2}{}\fi
\noindent\textit{Source reference.} Original reference 05W1: focused-seeds/geometry/S9-spaces-unramified-univ-inj/ref-tags.txt, source lines 4--4. Navigation pointer only: the counted proof cites this statement by \textbackslash{}ref as a standing input; the deduction the counted proof has to make is carried out inside the counted span, and no part of that deduction is delegated to this pointer.\par\medskip
\ifdefined\hypertarget\hypertarget{sourcepointer-3}{}\fi
\noindent\textit{Source reference.} Original reference 040X: focused-seeds/geometry/S9-spaces-unramified-univ-inj/ref-tags.txt, source lines 5--5. Navigation pointer only: the counted proof cites this statement by \textbackslash{}ref as a standing input; the deduction the counted proof has to make is carried out inside the counted span, and no part of that deduction is delegated to this pointer.\par\medskip
\ifdefined\hypertarget\hypertarget{sourcepointer-4}{}\fi
\noindent\textit{Source reference.} Original reference 03IE: focused-seeds/geometry/S9-spaces-unramified-univ-inj/ref-tags.txt, source lines 6--6. Navigation pointer only: the counted proof cites this statement by \textbackslash{}ref as a standing input; the deduction the counted proof has to make is carried out inside the counted span, and no part of that deduction is delegated to this pointer.\par\medskip
\ifdefined\hypertarget\hypertarget{sourcepointer-5}{}\fi
\noindent\textit{Source reference.} Original reference 03BZ: focused-seeds/geometry/S9-spaces-unramified-univ-inj/ref-tags.txt, source lines 7--7. Navigation pointer only: the counted proof cites this statement by \textbackslash{}ref as a standing input; the deduction the counted proof has to make is carried out inside the counted span, and no part of that deduction is delegated to this pointer.\par\medskip
\ifdefined\hypertarget\hypertarget{sourcepointer-6}{}\fi
\noindent\textit{Source reference.} Original reference 03H4: focused-seeds/geometry/S9-spaces-unramified-univ-inj/ref-tags.txt, source lines 8--8. Navigation pointer only: the counted proof cites this statement by \textbackslash{}ref as a standing input; the deduction the counted proof has to make is carried out inside the counted span, and no part of that deduction is delegated to this pointer.\par\medskip
\begin{lemma}
\label{lemma-where-unramified-universally-injective}
Consider a commutative diagram
$$
\xymatrix{
X \ar[rr]_f \ar[rd]_p & & Y \ar[ld]^q \\
& Z
}
$$
of algebraic spaces. Assume that
\begin{enumerate}
\item $p$ is locally of finite type,
\item $p$ is closed, and
\item $p_2 : X \times_Y X \to Z$ is closed.
\end{enumerate}
Then there exists an open subspace $W \subset Z$
such that a morphism $Z' \to Z$ factors through $W$ if and only if the
base change $f_{Z'} : X_{Z'} \to Y_{Z'}$ is unramified and universally
injective.
\end{lemma}

\begin{proof}
After replacing $Z$ by the open subspace found in
Lemma \ref{lemma-where-unramified}
we may assume that $f$ is already unramified; note that this does not
destroy assumption (2) or (3). By
Morphisms of Spaces, Lemma
\ref{spaces-morphisms-lemma-diagonal-unramified-morphism}
we see that $\Delta_{X/Y} : X \to X \times_Y X$ is an open immersion.
This remains true after any base change. Hence by
Morphisms of Spaces, Lemma
\ref{spaces-morphisms-lemma-universally-injective}
we see that $f_{Z'}$ is universally injective if and only if
the base change of the diagonal $X_{Z'} \to (X \times_Y X)_{Z'}$
is an isomorphism. Let $W \subset Z$ be the open subspace
(see
Properties of Spaces, Lemma
\ref{spaces-properties-lemma-open-subspaces})
with underlying set of points
$$
|W| = |Z| \setminus
|p_2|\left(|X \times_Y X| \setminus \Im(|\Delta_{X/Y}|)\right)
$$
i.e., $z \in |Z|$ is a point of $W$ if and only if the fibre of
$|X \times_Y X| \to |Z|$ over $z$ is in the image of
$|X| \to |X \times_Y X|$. Then it is clear from the discussion above
that the restriction $p^{-1}(W) \to q^{-1}(W)$ of $f$ is
unramified and universally injective.

\medskip\noindent
Conversely, suppose that $f_{Z'}$ is unramified and universally injective.
In order to show that $Z' \to Z$ factors through $W$ it suffices to show
that $|Z'| \to |Z|$ has image contained in $|W|$, see
Properties of Spaces, Lemma
\ref{spaces-properties-lemma-factor-through-open-subspace}.
Hence it suffices to prove the result when $Z'$ is the spectrum of a field.
Denote $z \in |Z|$ the image of $|Z'| \to |Z|$. The discussion above shows
that
$$
|X_{Z'}| \longrightarrow |(X \times_Y X)_{Z'}|
$$
is surjective. By
Properties of Spaces,
Lemma \ref{spaces-properties-lemma-points-cartesian}
in the commutative diagram
$$
\xymatrix{
|X_{Z'}| \ar[d] \ar[r] &
|(X \times_Y X)_{Z'}| \ar[d] \\
|p|^{-1}(\{z\}) \ar[r] &
|p_2|^{-1}(\{z\})
}
$$
the vertical arrows are surjective. It follows that $z \in |W|$ as desired.
\end{proof}

\end{document}
